
%=============================================================================
% Appendix O (CLEAN, v38-compatible): α^{57} exponent from primed determinant scaling
% and the explicit observer-dictionary step to the measured invariant I = G ħ H_0^2 / c^5.
%
% IMPORTANT: This appendix is theorem-grade on the *mathematical* side (finite-mode
% determinant scaling). The single remaining physics step is stated explicitly as
% a dictionary postulate (Observer Dictionary), consistent with v38's own language
% that Eq.(384) was previously "proposed".
%=============================================================================

\section{The $\alpha^{57}$ Mode-Count Exponent and the $G$--$H_0$--$\alpha$ Dictionary}
\label{app:O_alpha57_clean}

\subsection{O.1\quad Mathematical core: primed-determinant scaling fixes the exponent}
Let $\mathcal{H}$ be a finite-dimensional complex Hilbert space of dimension $k_{\max}$, and let
$\mathcal{K}:\mathcal{H}\to\mathcal{H}$ be a self-adjoint, positive semidefinite operator with
$\dim\ker(\mathcal{K})=N_{\rm gen}$.
Denote by $\det'(\mathcal{K})$ the \emph{primed determinant} over the nonzero spectrum of $\mathcal{K}$.

\begin{lemma}[Primed determinant scaling]
\label{lem:O_det_scaling}
For any $g>0$,
\begin{equation}
\det'\!\big(g\,\mathcal{K}\big)\;=\; g^{\,k_{\max}-N_{\rm gen}}\;\det'(\mathcal{K})\,.
\end{equation}
\end{lemma}

\begin{proof}
Diagonalize $\mathcal{K}$ on $\mathcal{H}$ with eigenvalues $\{\lambda_i\}_{i=1}^{k_{\max}}$.
Exactly $N_{\rm gen}$ of these are zero; the remaining $N:=k_{\max}-N_{\rm gen}$ satisfy $\lambda_i>0$.
Then by definition $\det'(\mathcal{K})=\prod_{i=1}^{N}\lambda_i$ (product over the nonzero spectrum), and
$\det'(g\mathcal{K})=\prod_{i=1}^{N}(g\lambda_i)=g^{N}\prod_{i=1}^{N}\lambda_i$.
\end{proof}

\begin{definition}[Microsector hierarchy factor as a \emph{determinant ratio}]
\label{def:O_eps}
Define
\begin{equation}
\varepsilon(g)\;:=\;\frac{\det'(\mathcal{K})}{\det'(g\,\mathcal{K})}\,.
\end{equation}
\end{definition}

\begin{corollary}[Topologically forced exponent]
\label{cor:O_alpha57}
If $k_{\max}=60$ and $N_{\rm gen}=3$, then
\begin{equation}
\varepsilon(g)=g^{-57}\,,
\qquad\text{and in particular}\qquad
\varepsilon(\alpha^{-1})=\alpha^{57}\,.
\end{equation}
\end{corollary}

\begin{proof}
Immediate from Lemma~\ref{lem:O_det_scaling} and Definition~\ref{def:O_eps} with $N=k_{\max}-N_{\rm gen}=57$.
\end{proof}

\subsection{O.2\quad Why this is \emph{not} a partition function (and why that matters)}
The object $\varepsilon(g)$ is a \emph{ratio of primed determinants} and therefore a \emph{dimensionless
coefficient ratio}. No claim is made that it equals a thermodynamic partition function, and no use is
made of $\log Z$.

A convenient way to interpret $\varepsilon(g)$ is as the ratio of Gaussian normalization constants for a
finite-dimensional quadratic form on the nonzero-mode quotient $\mathcal{H}/\ker(\mathcal{K})$:
if one defines for $g>0$
\begin{equation}
\mathcal{N}(g)\;:=\;\int_{\mathbb{C}^{N}} \exp\!\Big(-\,\langle \phi,\,(g\mathcal{K}_+)\phi\rangle\Big)\,d^{2N}\phi
\;\propto\; \frac{1}{\det(g\mathcal{K}_+)}\,,
\end{equation}
then $\mathcal{N}(1)/\mathcal{N}(g)=\det(\mathcal{K}_+)/\det(g\mathcal{K}_+)=\varepsilon(g)$, where
$\mathcal{K}_+$ is $\mathcal{K}$ restricted to the nonzero spectrum.
This is standard finite-dimensional Gaussian integration and fixes the \emph{power of $g$} to be the
nonzero-mode count $N$ (here $N=57$).

\subsection{O.3\quad The observer dictionary step (explicit)}
Define the observed dimensionless invariant
\begin{equation}
I\;:=\;\frac{G\,\hbar\,H_0^2}{c^5}\,.
\end{equation}
As shown in the main text (critical density vs. Planck density algebra),
\begin{equation}
\frac{\rho_c}{\rho_{\rm Pl}}
\;=\;
\frac{3}{8\pi}\,I\,,
\qquad\text{and}\qquad
\frac{\rho_\Lambda}{\rho_{\rm Pl}}
\;=\;
\Omega_\Lambda\,\frac{3}{8\pi}\,I\,.
\label{eq:O_dictionary_rho}
\end{equation}

\begin{definition}[Dictionary postulate: identifying the hierarchy factor]
\label{def:O_dictionary}
In the non-extensive microsector picture, the (dimensionless) hierarchy between the
cosmological IR scale and the Planck scale is encoded by the microsector coefficient ratio
$\varepsilon(\alpha^{-1})$. The observer dictionary identifies this hierarchy with the measured
invariant $I$:
\begin{equation}
\boxed{\;I\;=\;\varepsilon(\alpha^{-1})\;}\,.
\label{eq:O_dictionary_I}
\end{equation}
\end{definition}

\begin{theorem}[$G$--$H_0$--$\alpha$ invariant (dictionary-closed)]
\label{thm:O_I_equals_alpha57}
Assume $k_{\max}=60$ and $N_{\rm gen}=3$ (as derived in Appendix~\ref{app:gauge-emergence}), and adopt the dictionary
identification~\eqref{eq:O_dictionary_I}. Then
\begin{equation}
\boxed{\;\frac{G\,\hbar\,H_0^2}{c^5}\;=\;\alpha^{57}\;}\,.
\end{equation}
Consequently,
\begin{equation}
\frac{\rho_c}{\rho_{\rm Pl}}=\frac{3}{8\pi}\,\alpha^{57},
\qquad
\frac{\rho_\Lambda}{\rho_{\rm Pl}}=\Omega_\Lambda\,\frac{3}{8\pi}\,\alpha^{57}.
\end{equation}
\end{theorem}

\begin{proof}
By Corollary~\ref{cor:O_alpha57}, $\varepsilon(\alpha^{-1})=\alpha^{57}$. The theorem then follows
immediately from the dictionary identification~\eqref{eq:O_dictionary_I} and the algebraic relations
\eqref{eq:O_dictionary_rho}.
\end{proof}

\subsection{O.4\quad Status and integration guidance}
\begin{itemize}
\item Lemma~\ref{lem:O_det_scaling} and Corollary~\ref{cor:O_alpha57} are \emph{pure mathematics} and are theorem-grade.
\item The remaining step is the explicit observer-dictionary identification (Definition~\ref{def:O_dictionary}),
which upgrades the previously proposed relation into a \emph{closed dictionary statement} without any appeal to $\log Z$.
\end{itemize}
